% Contenido del póster: Lazos (2025), Technological Forecasting and Social Change. DOI 10.1016/j.techfore.2025.124119
% Elaborado con apoyo de IA a partir del resumen oficial del artículo (OpenAlex), 2026-11-12.
\def\DOI{10.1016/j.techfore.2025.124119}
\def\TituloEs{Cuando los reguladores quieren proteger a los inversionistas del crowdfunding de capital: ¿tiene efectos regionales?}
\def\TituloOriginal{When policy makers want to protect investors in equity crowdfunding: Does it have regional implications?}
\def\Autores{Lazos}
\def\Anio{2025}
\def\Revista{Technological Forecasting and Social Change}

\def\Problema{Las reglas que protegen a los inversionistas pueden cambiar qué empresas consiguen financiamiento, incluidas las de regiones que suelen recibir menos capital.}
\def\Objetivo{Investigar si una política de protección a los inversionistas en empresas jóvenes no cotizadas beneficia a las empresas de regiones con menos acceso al financiamiento con capital.}
\def\Metodo{Estudio del mercado de crowdfunding de capital del Reino Unido, donde se introdujo un tope de inversión solo para los inversionistas no acreditados. Se compara la probabilidad de éxito de las empresas de regiones desfavorecidas.}
\def\Resultados{\begin{itemize}
  \item El tope de inversión para inversionistas no acreditados puede aumentar la probabilidad de que las empresas de regiones desfavorecidas consigan capital.
  \item Aun así, esas empresas siguen teniendo menos probabilidad de completar con éxito su emisión.
\end{itemize}}
\def\Conclusion{Una regulación pensada para proteger al inversionista puede tener efectos distributivos regionales favorables, aunque no elimina la desventaja de las regiones con menos acceso al capital.}
\def\Relevancia{La modificación peruana de noviembre de 2025 amplió los límites de financiamiento por proyecto y por receptor (hasta 500 UIT por proyecto empresarial, o 1\,000 UIT con criterios de selección específicos), no los de inversión. Este estudio sugiere evaluar no solo cuánto se financia tras un cambio normativo, sino a quién se financia.}
\def\PalabrasClave{crowdfunding de capital, protección al inversionista, topes de inversión, regiones, Reino Unido}
% --- Ficha del estudio (herramientas/posteres/ficha_articulos.py) ---
\def\Ambito{Reino Unido (crowdfunding de capital)}
\def\Enfoque{Empírico: efecto de un tope de inversión}
\def\FiguraPeru{1}
% --- fin de la ficha ---
